\begin{lemma}[The spectral projection removes exactly the Gaussian mean]
Write $\mathfrak{m}(v) = \int v \, d\gamma_n$ for $v \in L^2(\gamma_n)$.  Then for
every $v \in L^2(\gamma_n)$,
\[
  \Pi_0 v = v - \mathfrak{m}(v)\, h_0
  \qquad\text{in } L^2(\gamma_n),
\]
where $\Pi_0 = L^0$ is the operator of \noderef{FracPow} and $h_0 \in L^2(\gamma_n)$
is the Hermite element of index $0$ from \noderef{HermiteL2}.  Since $h_0$ is the
constant function $1$ by \noderef{Hermite1Zero}, this is the identity
$\Pi_0 v = v - \mathfrak{m}(v)\mathbf{1}$ of the reference paper.
\end{lemma}
\begin{proof}
Both sides are elements of $L^2(\gamma_n)$.  By \noderef{HermiteBasis} the Hermite
coefficient map $w \mapsto (w_\beta)_\beta$, $w_\beta = \langle w, h_\beta\rangle$,
is an injective linear isometry onto $\ell^2(\mathbb{N}_0^n)$; so it suffices to
check that the two sides have the same coefficient at every
$\beta \in \mathbb{N}_0^n$.

\emph{Step 1: the zeroth coefficient of $v$ is its mean.}
By \noderef{Hermite1Zero} the function $h_0$ is the constant $1$, and
$\gamma_n$ is a probability measure by \noderef{GaussIsProb}.  Hence
\[
  v_0 = \langle v, h_0\rangle = \int v \cdot 1 \, d\gamma_n = \mathfrak{m}(v),
\]
the integral being finite because $v \in L^2(\gamma_n) \subseteq L^1(\gamma_n)$ by
\noderef{L2IntoL1}.

\emph{Step 2: the coefficients of the right-hand side.}
By \noderef{HermiteL2Orthonormal}, $\langle h_0, h_\beta\rangle = 1$ for
$\beta = 0$ and $= 0$ for $\beta \neq 0$.  By linearity of the coefficient map,
the element $v - \mathfrak{m}(v) h_0$ therefore has coefficients
\[
  \bigl(v - \mathfrak{m}(v)h_0\bigr)_0 = v_0 - \mathfrak{m}(v) = 0,
  \qquad
  \bigl(v - \mathfrak{m}(v)h_0\bigr)_\beta = v_\beta \quad (\beta \neq 0),
\]
using Step 1 for the first equality.

\emph{Step 3: the coefficients of the left-hand side.}
By \noderef{FracPowCoeff} with $s = 0$, the multiplied coefficient family is
\[
  m_0\bigl((v_\beta)_\beta\bigr)_\beta
    = \begin{cases} 0, & \beta = 0,\\ |\beta|^0 v_\beta = v_\beta, & \beta \neq 0.\end{cases}
\]
This family is square summable, being dominated termwise by $(|v_\beta|^2)_\beta$,
which is summable with sum $\|v\|_2^2$ by Parseval (\noderef{HermiteBasis}).  By
\noderef{FracPowDomain} with $s = 0$ this says $v \in \operatorname{Dom}(L^0)$, so
the defining formula of \noderef{FracPow} applies and gives $\Pi_0 v$ precisely
those coefficients: $0$ at $\beta = 0$ and $v_\beta$ at $\beta \neq 0$.

Steps 2 and 3 produce the same coefficient family, so the two elements of
$L^2(\gamma_n)$ coincide.
\end{proof}
